% ===== Block B02: 原文 p.17–25 | §1.2 稳定 $N$-点曲线 | 译手 tr-ch1b =====
\subsection{稳定 $N$-点曲线}
% [原文 p.17]
\noindent\textbf{1.2.1 结点曲线.}\quad 所谓结点曲线（nodal curve）$C$，是指复射影空间 $\PP^n$ 中的一个闭解析子簇，它的奇异点都是结点，也就是说，它在局部解析同构于 $zw=0$ 在原点处的邻域 $U_\eps=\{(z,w)\mid zw=0,\ |z|<\eps,\ |w|<\eps\}$。若使用代数几何的语言，则结点曲线就是半稳定曲线（semi-stable curve），即定义在复数域 $\C$ 上的既约连通完备代数曲线。

设 $P$ 是结点曲线 $C$ 的一个结点。当 $\eps$ 充分小时，$P$ 有一个邻域同构于上述 $U_\eps$。于是，把 $U_+=U_\eps\cap\{w=0\}=\{z\mid |z|<\eps\}$ 与 $U_-=U_\eps\cap\{z=0\}=\{w\mid |w|<\eps\}$ 分离开来，就可以消解结点 $P$ 处的奇点。这样，代替一个点 $P$，我们得到两个点 $P_+=0\in U_+$ 与 $P_-=0\in U_-$。如此得到的新簇 $\tilde C$ 称为在点 $P$ 处作正规化所得的曲线。存在自然的全纯映射 $\varpi:\tilde C\to C$，使得 $\varpi$ 是从 $\tilde C\setminus\{P_+,P_-\}$ 到 $C\setminus\{P\}$ 的恒等映射，并且 $\varpi^{-1}(P)=\{P_+,P_-\}$。

% [Tutor Note: 术语首次出现：glossary 给 nodal curve => 结点曲线（本块统一用此译名），
%  semi-stable curve => 半稳定曲线，node => 结点；$\varpi$ 是原文的软体 varpi，
%  不是 $\pi$——§1.3 里同一个正规化映射换成 $\pi$ 或 $\nu$，读时勿混。
%  "分开 $U_+$ 与 $U_-$ 即消解结点"的实质：局部环 $\C\{z,w\}/(zw)$ 的两个分支
%  $z=0$ 与 $w=0$ 各对应一个光滑分支，正规化把这个奇点"拆成"两个非奇异点
%  $P_+,P_-$；核对：$U_+\cap U_-=\{(0,0)\}$ 原本是 $C$ 上唯一结点 $P$，分开后成为
%  两个点，故 $\varpi$ 在 $P_\pm$ 之外是双全纯。原文 p.17]

反之，从紧黎曼曲面出发也可以构造结点曲线。设 $P_+$ 与 $P_-$ 是紧黎曼曲面 $R$ 上两个不同的点。取 $R$ 上以 $P_+$（resp. $P_-$）为中心的局部坐标 $z$（resp. $w$）。我们把 $P_+$ 与 $P_-$ 认同，由此得到一个新空间 $C$。用 $P$ 表示 $C$ 中对应于 $P_\pm$ 的点。我们把 $C$ 中 $P$ 的一个邻域与 $zw=0$ 在原点处的邻域认同，并把 $C\setminus\{P\}$ 与 $R\setminus\{P_+,P_-\}$ 认同，则 $C$ 是结点曲线。若在 $C\setminus\{P\}$ 上再取另外两个点并重复同样的过程，就得到有两个结点的结点曲线。这样，任何结点曲线都可以由紧黎曼曲面得到。

% [原文 p.18]
\figomit{1.2}{亏格 0 的稳定 3-点曲线。}

\noindent\textbf{1.2.2 稳定 $N$-点曲线.}\quad

\begin{definition}[1.16]
由一条曲线 $C$ 以及 $C$ 上的点 $Q_1,Q_2,\dots,Q_N$ 组成的数据 $\mathfrak{X}=(C;Q_1,Q_2,\dots,Q_N)$ 称为稳定 $N$-点曲线（stable $N$-pointed curve），若下列条件被满足。
\begin{enumerate}[label=(\arabic*)]
\item 曲线 $C$ 是紧黎曼曲面或结点曲线。
\item $Q_1,Q_2,\dots,Q_N$ 是曲线 $C$ 的非奇异点。
\item 若不可约分量 $C_i$ 是射影直线（即黎曼球面）$\PP^1$（resp. 带一个二重点的有理曲线，resp. 椭圆曲线），则 $C_i$ 与其他分量的交点个数再加上 $C_i$ 上 $Q_j$ 的个数，至少为 $3$（resp. $1$）。
\item $\dim_{\C}H^1(C,\cO_C)=g$。
\end{enumerate}
\end{definition}

% [Tutor Note: 条件 (3) 的三情形只给两个阈值，须自己对上号：$\PP^1$ 与"带一个二重点
%  的有理曲线"都要求 $\ge 3$，椭圆曲线（含带一个结点的椭圆曲线）要求 $\ge 1$。
%  核对：亏格 0 时 $C_i=\PP^1$ 的自同构群 3 维可迁作用于 3 点，故 $N\ge3$（例 1.17）；
%  亏格 1 时平移 1 维，故至少要 1 个点才能杀死平移，这正是阈值为 1 的来源。原文 p.18]
% [Tutor Note: 条件 (4) 用的是 $H^1(C,\cO_C)$ 的维数而非"算术亏格"：对结点曲线
%  $\dim H^1(C,\cO_C)=g$ 正是几何亏格的定义（对偶化层 $\omega_C$ 的次数为 $2g-2$）；
%  后文 (1.50) 一带的 Serre 对偶 $H^1(C,\cO_C(-m\Sigma Q_j))^\ast\simeq H^0(C,\omega_C(m\Sigma Q_j))$
%  就用这条。原文 p.18]

注意，上述条件 (3) 等价于说 $\Aut(\mathfrak{X})$ 是有限群，即 $\mathfrak{X}$ 没有无穷小自同构。下文我们常常对 $N$-点稳定曲线 $\mathfrak{X}$ 再加下述条件 (Q)。

(Q)\quad 每个分量 $C_i$ 至少含有一个 $Q_j$。

条件 (Q) 的含义将在下面的引理 1.21 与引理 1.22 中说明。

\begin{example}[1.17]
对 $g=0$，由性质 (3) 得 $N\ge 3$。亏格 0 的任何 3-点稳定曲线都同构于图 1.2 中的曲线。

亏格 0 的任何 4-点稳定曲线都同构于图 1.3 中下述两条稳定曲线之一。

亏格 1 的任何 1-点稳定曲线都同构于图 1.4 中下述两条稳定曲线之一。
\end{example}

% [Tutor Note: 条件 (3) 与"$\Aut(\mathfrak{X})$ 有限"的等价性要按分量读：$\PP^1$ 上
%  $\le2$ 个标记点时仍有 $\C^\ast$ 或 $\C_+$ 型的单参数自同构，带一个二重点的有理曲线
%  没有标记点时可以沿结点"旋转"；两条都违反"有限"。反之若每条 $\PP^1$ 分量上
%  (交点数 + 标记数) $\ge3$、椭圆分量上 $\ge1$，则自同构被钉死。原文 p.18]

\begin{definition}[1.18]
设 $C$ 是曲线，$Q$ 是 $C$ 上的非奇异点。$C$ 在点 $Q$ 处的第 $n$ 无穷小邻域（$n$-th infinitesimal neighbourhood）$s^{(n)}$ 是 $\C$-代数同构
\begin{equation}
s^{(n)}:\cO_{C,Q}/\mathfrak{m}_Q^{n+1}\simeq\C[[\xi]]/(\xi^{n+1}) \tag{1.42}
\end{equation}
其中 $\mathfrak{m}_Q$ 是 $\cO_{C,Q}$ 的极大理想，由在 $Q$ 处消失的全纯函数芽组成。
\end{definition}

% [原文 p.19]
\figomit{1.3}{亏格 0 的稳定 4-点曲线。}
\figomit{1.4}{亏格 1 的稳定 1-点曲线。}

% [Tutor Note: 图 1.2 是 $\PP^1=\C\cup\{\infty\}$ 上标出 $0,1,\infty$ 三点，它给出
%  亏格 0 的 3-点稳定曲线的唯一同构类（见例 1.17）；三点的"位置"无意义，因为
%  $PGL(2,\C)$ 可迁地作用在有序三点组上。图 1.3 左：一条 $\PP^1$ 上四个标记点；
%  右：两条 $\PP^1$ 在结点处相接，每支上两个标记点。图 1.4 左：光滑椭圆曲线带一点 $Q$；
%  右：带一个结点的椭圆曲线带一点 $Q$。后两图就是例 1.17 所说的"各两种"，分别是
%  $g=0,N=4$ 与 $g=1,N=1$ 稳定曲线的全部同构类（模空间的内点与边界点各一个）。原文 p.18–p.19]

在 (1.42) 中令 $n\to\infty$，便得到同构
\begin{equation}
s^{(\infty)}:\hat\cO_{C,Q}\simeq\C[[\xi]]. \tag{1.43}
\end{equation}
同构 $s^{(\infty)}$ 称为 $C$ 在 $Q$ 处的形式邻域（formal neighbourhood）。

% [Tutor Note: (1.43) 只是把 (1.42) 取逆极限：$\hat\cO_{C,Q}=\varprojlim_n\cO_{C,Q}/\mathfrak{m}_Q^{n+1}$，
%  而 $\varprojlim_n\C[[\xi]]/(\xi^{n+1})=\C[[\xi]]$。注意 (1.43) 的靶空间没有分母，
%  即形式邻域是"无穷阶"的坐标，而 (1.42) 是它 mod $\xi^{n+1}$ 的截断；定义 1.19 与
%  定义 1.23 后来正是用这个无穷阶坐标定义带形式邻域的族。原文 p.19]

\begin{definition}[1.19]
数据 $\mathfrak{X}=(C;Q_1,Q_2,\dots,Q_N;s_1,s_2,\dots,s_N)$ 称为亏格 $g$ 的带形式邻域的稳定 $N$-点曲线（stable $N$-pointed curve of genus $g$ with formal neighbourhoods），若下列条件被满足。
\begin{enumerate}[label=(\arabic*)]
\item $(C;Q_1,Q_2,\dots,Q_N)$ 是亏格 $g$ 的 $N$-点稳定曲线。
\item $s_j$ 是 $C$ 在 $Q_j$ 处的形式邻域。
\end{enumerate}
\end{definition}

类似地，我们定义带第 $n$ 形式邻域的稳定 $N$-点曲线
\[
\mathfrak{X}^{(n)}=(C;Q_1,Q_2,\dots,Q_N;s_1^{(n)},s_2^{(n)},\dots,s_N^{(n)}).
\]

\noindent\textbf{1.2.3 留数配对.}\quad 设 $C$ 是结点曲线，即 $C$ 是至多只有通常二重点的既约曲线，并且在 $\C$ 上固有。设 $\Omega^1_C$ 是 $C$ 的 Kähler 微分层，$\omega_C$ 是 $C$ 的对偶化层。在奇异点 $P$ 附近，$C$ 解析同构于由 $xy=0$ 定义的簇。在这组坐标下，层 $\Omega^1_C$ 表为
\begin{equation}
\Omega^1_C=(\cO_C\dd x+\cO_C\dd y)/(x\,\dd y+y\,\dd x)\cO_C. \tag{1.44}
\end{equation}
另一方面，在奇异点 $P$ 附近，对偶化层 $\omega_C$ 是一个可逆层，它由微分 $\zeta$ 生成，而 $\zeta$ 在 $x\ne 0$ 处为 $\dd x/x$，在 $y\ne 0$ 处为 $-\dd y/y$。此外，在 $C$ 的奇异点之外，层 $\Omega^1_C$ 与 $\omega_C$ 重合。于是我们有
\begin{equation}
\Omega^1_C=\mathfrak{m}\,\omega_C, \tag{1.45}
\end{equation}

% [Tutor Note: (1.44) 的由来：$\Omega^1_{\C^2}=\cO\dd x\oplus\cO\dd y$ 模去 $xy=0$ 的关系
%  $\dd(xy)=y\,\dd x+x\,\dd y$。核对两个分支：在 $y=0$ 的支上关系化为 $x\,\dd y=0$，
%  故 $\Omega^1_C$ 由 $\dd x$ 生成（与 $U_+$ 上的 $\dd z$ 一致）；在 $x=0$ 的支上由 $\dd y$
%  生成。这说明 (1.44) 在结点处给出秩 1 的层，而不是秩 2 的自由层。原文 p.19]
% [Tutor Note: (1.45) 中 $\mathfrak{m}$ 是奇点集的定义理想层（见 p.20 首句），故 $\Omega^1_C$
%  在非奇异点与 $\omega_C$ 相同、在结点处严格小于 $\omega_C$；商恰是
%  $\omega_C\otimes\cO_{C_{\mathrm{sing}}}$，即 p.20 的正合列。留数配对之所以必须用
%  $\omega_C$（而不是 $\Omega^1_C$）：$\omega_C$ 在结点处允许一阶极点，其留数才有意义。原文 p.19]
% [原文 p.20]
其中 $\mathfrak{m}$ 是 $C$ 的奇点集的定义理想层。于是我们有下述正合列。
\[
0\longrightarrow\Omega^1_C\longrightarrow\omega_C\longrightarrow\omega_C\otimes\cO_{C_{\mathrm{sing}}}\longrightarrow 0 .
\]

设 $\nu:\tilde C\to C$ 是曲线 $C$ 的正规化。又设 $\{P_1,\dots,P_q\}$ 是曲线 $C$ 的二重点集，并对每个二重点 $P_i$ 置 $\nu^{-1}(P_i)=\{P_{i,+},P_{i,-}\}$。则我们有下述正合列。
\begin{equation}
0\longrightarrow\omega_C\longrightarrow\nu_\ast\omega_{\tilde C}\Bigl(\sum_i^k(P_{i,+}+P_{i,-})\Bigr)\xrightarrow{\ r\ }\bigoplus_i^q\C\longrightarrow 0, \tag{1.46}
\end{equation}
其中在每个二重点 $P_i$ 处，映射 $r$ 由
\[
\operatorname{res}_{P_{i,+}}(\tau)-\operatorname{res}_{P_{i,-}}(\tau)
\]
给出。

% [Tutor Note: (1.46) 的求和上限印作 $k$，但同一式中 $\bigoplus$ 的上限是 $q$，正文又
%  明确写 $\{P_1,\dots,P_q\}$ 是二重点集，故上限应为 $q$；$k$ 在全书中无定义。照原文排。
%  另：正合列左端的同构来源是 $\omega_C$ 的截面在 $P_{i,\pm}$ 处允许一阶极点且要求
%  $\operatorname{res}_{P_{i,+}}=\operatorname{res}_{P_{i,-}}$——这正是 $r$ 的核。
%  核对：对 $\C[[z,w]]/(zw)$，$\dd z/z+\dd w/w$ 型的截面给出 $r=0$，而
%  $\dd z/z-\dd w/w$ 型给出 $r\ne0$，故 $\omega_C$ 的截面是"留数和为零"的那一半。原文 p.20]

这说明对偶化层 $\omega_C$ 的局部全纯截面，可以看作 $\tilde C$ 上亚纯 1-形式的局部截面：它在 $P_{i,+}$ 与 $P_{i,-}$ 处各有一阶极点，留数之和为零，并且在 $P_{i,\pm}$ 之外全纯。下文我们常常使用这一解释。下述引理是这一解释的简单推论。

\begin{lemma}[1.20]
设 $\tau$ 是对偶化层 $\omega_C$ 的亚纯截面，且在二重点处全纯。则 $\tau$ 的留数之和为零。
\end{lemma}

\begin{lemma}[1.21]
设带形式邻域的 $N$-点稳定曲线满足条件 (Q)。用 $t_j$ 表示关于形式参数 $\xi_j=s_j^{-1}(\xi)$ 在 $Q_j$ 处的 Laurent 展开。则下述同态是单射：
\begin{equation}
t=\bigoplus t_j:H^0\Bigl(C,\cO\bigl(\ast\sum_{j=1}^N Q_j\bigr)\Bigr)\longrightarrow\bigoplus_{j=1}^N\C((\xi_j)), \tag{1.47}
\end{equation}
\begin{equation}
t=\bigoplus t_j:H^0\Bigl(C,\omega_C\bigl(\ast\sum_{j=1}^N Q_j\bigr)\Bigr)\longrightarrow\bigoplus_{j=1}^N\C((\xi_j))\dd\xi_j, \tag{1.48}
\end{equation}
其中 $\omega_C$ 是 $C$ 的对偶化层，并且
\[
H^0\Bigl(C,\cO\bigl(\ast\sum_{j=1}^N Q_j\bigr)\Bigr)=\varinjlim_{n\to\infty}H^0\Bigl(C,\cO\bigl(n\sum_{j=1}^N Q_j\bigr)\Bigr).
\]
\end{lemma}

% [Tutor Note: $s_j^{-1}(\xi)$ 的意思是：形式邻域 $s_j$ 把 $\cO_{C,Q_j}$ 的完备化
%  同构到 $\C[[\xi]]$，于是 $\xi$ 通过 $s_j^{-1}$ 成为 $C$ 在 $Q_j$ 附近的"形式坐标"，
%  $t_j$ 就是把亚纯函数按这个形式坐标展开。这一步是 §1.2 全部内容的技术核心：
%  条件 (Q) 保证展开映射单射，使函数/微分可以"只看有限个 $Q_j$ 处的 Laurent 尾"。原文 p.20]

\noindent\textit{证明.}\quad 由条件 (Q)，若 $f\in H^0(C,\cO(\ast\sum Q_j))$ 在某个不可约分量上为零，则 $f=0$，故 (1.47) 是单射。关于 $\omega$ 的情形同理可证。\hfill$\square$

由本引理，$H^0(C,\cO(\ast\sum Q_j))$（resp. $H^0(C,\omega_C(\ast\sum Q_j))$）可以看作 $\bigoplus_{j=1}^N\C((\xi_j))$（resp. $\bigoplus_{j=1}^N\C((\xi_j))\dd\xi_j$）的子空间。存在留数配对

% [原文 p.21]
\begin{equation}
\begin{gathered}
\bigoplus_{j=1}^N\C((\xi_j))\times\bigoplus_{j=1}^N\C((\xi_j))\dd\xi_j\longrightarrow\C,\\
\bigl((f(\xi_1),\dots,f(\xi_N)),(g(\xi_1)\dd\xi_1,\dots,g(\xi_N)\dd\xi_N)\bigr)
\longmapsto\sum_{j=1}^N\operatorname*{Res}_{\xi_j=0}\bigl(f(\xi_j)g(\xi_j)\dd\xi_j\bigr).
\end{gathered} \tag{1.49}
\end{equation}

下述定理在我们的理论中起重要作用。

\begin{theorem}[1.22]
在留数配对下，向量空间 $H^0(C,\cO(\ast\sum_{j=1}^N Q_j))$ 与向量空间 $H^0(C,\omega_C(\ast\sum_{j=1}^N Q_j))$ 互为对方的零化子。
\end{theorem}

\noindent\textit{证明.}\quad 由于留数与局部坐标的选取无关，我们可以假定形式邻域在点 $Q_j$ 的邻域内给出全纯坐标。对任意正整数 $m,n$，我们有正合列
\begin{equation}
0\longrightarrow\cO_C\Bigl(-m\sum_{j=1}^N Q_j\Bigr)\longrightarrow\cO_C\Bigl(n\sum_{j=1}^N Q_j\Bigr)\xrightarrow{\ p\ }\bigoplus_{j=1}^N\bigoplus_{k=-n}^{m-1}\C\xi_j^k\longrightarrow 0 . \tag{1.50}
\end{equation}
由这个正合列，我们得到上同调群的长正合列
\[
H^0\Bigl(C,\cO_C\bigl(n\sum_{j=1}^N Q_j\bigr)\Bigr)\longrightarrow\bigoplus_{j=1}^N\bigoplus_{k=-n}^{m-1}\C\xi_j^k\xrightarrow{\ c\ }H^1\Bigl(C,\cO_C\bigl(-m\sum_{j=1}^N Q_j\bigr)\Bigr)\longrightarrow 0 .
\]
注意，由 Serre 对偶，$H^1(C,\cO_C(-m\sum_{j=1}^N Q_j))$ 的对偶空间是
\[
H^0\Bigl(C,\omega_C\bigl(m\sum_{j=1}^N Q_j\bigr)\Bigr).
\]
在 $c$ 的像上，对偶配对
\[
\langle\ ,\ \rangle:H^1\Bigl(C,\cO_C\bigl(-m\sum_{j=1}^N Q_j\bigr)\Bigr)\times H^0\Bigl(C,\cO_C\bigl(m\sum_{j=1}^N Q_j\bigr)\Bigr)\longrightarrow\C
\]
由
\begin{equation}
\bigl\langle c\bigl((g_j(\xi_j))\bigr),\tau\bigr\rangle=\sum_{j=1}^N\operatorname*{Res}_{\xi_j=0}\bigl(g_j(\xi_j)\tau_j\bigr) \tag{1.51}
\end{equation}
给出，其中 $(g_j(\xi_j))\in\bigoplus_{j=1}^N\bigoplus_{k=-n}^{m-1}\C\xi_j^k$，$\tau\in H^0(C,\cO_C(m\sum_{j=1}^N Q_j))$，而 $\tau_j$ 是 $\tau$ 在点 $Q_j$ 处关于 $\xi_j$ 的 Laurent 展开。$(g_j(\xi_j))\in\bigoplus_{j=1}^N\bigoplus_{k=-n}^{m-1}\C\xi_j^k$ 是整体亚纯函数 $g\in H^0(C,\cO_C(-m\sum_{j=1}^N Q_j))$ 的像，当且仅当
\[
\sum_{j=1}^N\operatorname*{Res}_{\xi_j=0}\bigl(g_j(\xi_j)\tau_j\bigr)=0
\]
对任意 $\tau\in H^0(C,\omega_C(m\sum_{j=1}^N Q_j))$ 成立。(1.51) 的证明将在后面给出。

% [Tutor Note: (1.50) 的正合性核对：$p$ 取 $\xi_j$-展开中次数 $-n\le k\le m-1$ 的截断，
%  其核是"在 $Q_j$ 消没 $\ge m$ 阶"的截面，即 $\cO_C(-m\sum Q_j)$；右端的维数
%  $N(m+n)$ 与 Riemann–Roch 合起来给出 $c$ 的余核同构于 $H^1$。又 (1.51) 左端的
%  $c((g_j(\xi_j)))$ 是连接同态作用在 $\oplus\oplus\C\xi_j^k$ 的像上，故 $g_j$ 只含有限项
%  （$-n\le k\le m-1$），(1.51) 的求和是有限和。原文 p.21]
% [Tutor Note: 笔误：本段先说"$H^1(C,\cO_C(-m\sum Q_j))$ 的对偶是 $H^0(C,\omega_C(m\sum Q_j))$"，
%  紧接着却把配对写成含 $H^0(C,\cO_C(m\sum Q_j))$ 的形式，并在 (1.51) 后写
%  $\tau\in H^0(C,\cO_C(m\sum Q_j))$；按 Serre 对偶与下一段的"对任意
%  $\tau\in H^0(C,\omega_C(m\sum Q_j))$"，配对里的全纯函数层应为 $\omega_C(m\sum Q_j)$。
%  同一个笔误在 p.23 的 Dolbeault 配对处再现；照原文排，不改。原文 p.21]
% [原文 p.22]
设 $(f_j(\xi_j))\in\bigoplus_{j=1}^N\C((\xi_j))$ 经留数配对被 $H^0(C,\omega_C(\ast\sum Q_j))$ 零化。固定一个正整数 $m$。对
\[
f_j(\xi_j)=\sum_{k=-n_j}^{\infty}a_k^{(j)}\xi_j^k,\qquad j=1,2,\dots,N
\]
置
\[
f_{j,m}(\xi_j)=\sum_{k=-n_j}^{m-1}a_k^{(j)}\xi_j^k,\qquad j=1,2,\dots,N .
\]
置 $n=\max_j n_j$。由于 $(f_j(\xi_j))$ 被 $H^0(C,\omega_C(\ast\sum Q_j))$ 零化，$(f_{j,m}(\xi_j))$ 被 $H^0(C,\omega_C(m\sum_{j=1}^N Q_j))$ 零化。$(f_{j,m}(\xi_j))$ 是亚纯函数 $f^{(m)}\in H^0(C,\cO_C(n\sum_{j=1}^N Q_j))$ 的像。设 $f_j^{(m)}(\xi_j)$ 是 $f^{(m)}$ 在 $Q_j$ 处关于 $\xi_j$ 的 Laurent 展开。则
\[
g_j(\xi_j):=f_j(\xi_j)-f^{(m)}(\xi_j)\equiv 0\bmod \xi_j^m .
\]
由于 $f^{(m)}$ 是整体亚纯函数，$(f_j^{(m)}(\xi_j))$ 被
\[
H^0\Bigl(C,\omega_C\bigl(\ast\sum_{j=1}^N Q_j\bigr)\Bigr)
\]
零化。于是 $(g_j(\xi_j))$ 也被 $H^0(C,\omega_C(\ast\sum Q_j))$ 零化。设 $g_k(\xi_k)\ne 0$，于是
\[
g_k(\xi_k)=b_s\xi_k^s+b_{s+1}\xi_k^{s+1}+\cdots,\qquad b_s\ne 0 .
\]
则 $s\ge m>0$。因此存在亚纯 1-形式
\[
\omega\in H^0(C,\omega_C((s+1)Q_k)),
\]
它在 $Q_k$ 处有 $s+1$ 阶极点，而在 $C\setminus Q_k$ 上全纯。于是我们有
\[
\sum_{j=1}^N\operatorname*{Res}_{\xi_j=0}g_j\omega_j=\operatorname*{Res}_{\xi_k=0}g_k\omega_k=b_s\ne 0,
\]
其中 $\omega_j$ 是 $\omega$ 在 $Q_j$ 处关于 $\xi_j$ 的 Laurent 展开。这与 $\{g_j(\xi_j)\}$ 被 $H^0(C,\omega_C(m\sum_{j=1}^N Q_j))$ 零化矛盾。于是我们断言 $g_j(\xi_j)=0$，$j=1,2,\dots,N$。因此 $f_j(\xi_j)$ 就是 $f^{(m)}$ 在 $Q_j$ 处关于 $\xi_j$ 的 Laurent 展开。

% [Tutor Note: 记号不一致：原文先写 $g_j(\xi_j):=f_j(\xi_j)-f^{(m)}(\xi_j)$（第二项无下标 $j$），
%  下一行又写 $(f_j^{(m)}(\xi_j))$；按上下文 $f^{(m)}(\xi_j)$ 应读作 $f_j^{(m)}(\xi_j)$，
%  即 $f^{(m)}$ 在 $Q_j$ 处的展开。照原文排，不改。原文 p.22]
% [Tutor Note: 关键跳步（★）：为什么存在 $\omega\in H^0(C,\omega_C((s+1)Q_k))$ 在 $Q_k$
%  有 $s+1$ 阶极点、别处全纯？用 Riemann–Roch 于 $K+D$（$D=(s+1)Q_k$）：
%  $h^0(K+D)-h^0(-D)=\deg(K+D)-g+1=g-1+s+1=s+g$，而 $h^0(-D)=0$（$s+1>0$），
%  故 $h^0(K+D)=g+s>g=h^0(\omega_C)$，超出的部分正是有极点的形式。原文 p.22]
% [Tutor Note: 截断为什么合法：$H^0(C,\omega_C(m\sum Q_j))$ 中的 $\tau$ 在 $Q_j$ 处极点阶
%  $\le m$，故 $\operatorname{Res}(g_j\tau_j)$ 只看见 $g_j$ 的 $k\le m-1$ 项，
%  $(f_j)$ 与 $(f_{j,m})$ 给出同一配对值；这就是"零化 $\Rightarrow$ 截断也零化"的来源。原文 p.22]

下证 (1.51)。先设 $C$ 是紧黎曼曲面。对小正数 $\eps$ 置 $D_\eps^{(j)}=\{\xi_j\mid |\xi_j|<\eps\}$ 以及 $C_\eps=C\setminus\bigcup_{j=1}^N D_\eps^{(j)}$。沿一条闭曲线切开 $C_\eps$，使所得的切面 $R'_\eps$ 是单连通的（见图 1.5）。

选取 $C$ 的一个开覆盖 $\{U_i\}_{i=1}^t$，使得
\[
Q_j\in U_j,\quad j=1,2,\dots,N,\qquad Q_j\notin U_i,\quad i\ne j .
\]
取 $U_i$ 足够小，使得对 $j=1,\dots,N$ 存在
\[
h_j\in H^0\Bigl(U_j,\cO_C\bigl(m\sum_{j=1}^N Q_j\bigr)\Bigr)
\]

% [原文 p.23]
\figomit{1.5}{沿闭曲线切开曲面 $C_\eps$。}
且 $p(h_j)=g_j(\xi_j)$。置 $h_i=0$（$i\ge N+1$）。则
\[
h_{ab}=h_b-h_a\in H^0\Bigl(U_a\cap U_b,\cO_C\bigl(-n\sum_{j=1}^N Q_j\bigr)\Bigr),
\]
并且 $\{h_{ab}\}$ 是一个上链，它代表
\[
c((g_j(\xi_j))\in H^1\Bigl(C,\cO_C\bigl(-n\sum_{j=1}^N Q_j\bigr)\Bigr).
\]
对偶配对 $H^1(C,\cO_C(-m\sum_{j=1}^N Q_j))\times H^0(C,\cO_C(n\sum_{j=1}^N Q_j))\to\C$ 由积分
\[
\frac{1}{2\pi\ii}\int_C\omega\wedge\tau
\]
给出，其中 $\omega\in H^{(0,1)}_{\bar\partial}(C,\cO_C(-m\sum_{j=1}^N Q_j))$ 是一个 Dolbeault 上同调类。设 $\cD_C^{(0,0)}$（resp. $\cD_C^{(0,k)}$）是 $C$ 上复值 $C^\infty$ 函数（resp. $(0,k)$ 型形式）芽的层，并定义
\[
\cD_C^{(0,i)}\Bigl(-n\sum_{j=1}^N Q_j\Bigr)
=\cO_C\Bigl(-n\sum_{j=1}^N Q_j\Bigr)\otimes_{\cO_C}\cD_C^{(0,i)} .
\]
于是我们有
\[
H^{(0,1)}_{\bar\partial}\Bigl(C,\cO_C\bigl(-m\sum_{j=1}^N Q_j\bigr)\Bigr)
\simeq\frac{\splitfrac{\operatorname{Ker}\Bigl\{\bar\partial:H^0\bigl(C,\cD_C^{(0,1)}(-m\sum_{j=1}^N Q_j)\bigr)\to}
{H^0\bigl(C,\cD_C^{(0,2)}(-m\sum_{j=1}^N Q_j)\bigr)\Bigr\}}}
{\bar\partial\Bigl(H^0\bigl(C,\cD_C^{(0,0)}(-m\sum_{j=1}^N Q_j)\bigr)\Bigr)} .
\]

% [Tutor Note: 图 1.5 里切开的目的：$\partial C_\eps=\bigcup_j\partial D_\eps^{(j)}$ 只有小圆周，
%  但要把它写成 $-\partial$ 的积分须保证 $C_\eps$ 上 Stokes 可用；沿 $\alpha_i,\beta_i$
%  与割线 $l_j$ 的积分两两相消（这正是下一段括注的来源），剩下的都是小圆周。原文 p.23]

% [原文 p.24]
\begin{sloppypar}
与 $\{h_{ab}\}=c(g_j(\xi_j))$ 相伴的 Dolbeault 上同调类如下给出。存在 $s_a\in H^0\bigl(U_a,\cD_C^{(0,0)}(-n\sum_{j=1}^N Q_j)\bigr)$ 使得
\end{sloppypar}
\[
h_{ab}=s_b-s_a .
\]
于是对任意 $a,b$ 我们有
\[
h_a-s_a=h_b-s_a .
\]
因此它们定义了 $C$ 上的一个整体函数 $h$。注意 $h$ 在 $Q_j$ 处有极点，故 $\bar\partial h$ 不是零上同调类。类 $\bar\partial h$ 代表与 Čech 一余循环 $\{h_{ab}\}$ 相伴的 Dolbeault 上同调类：
\[
\bar\partial h=-\bar\partial s_a
\]
在 $U_a$ 上。于是配对 $\langle c((g_j(\xi_j)),\tau\rangle$ 计算如下。
\[
\begin{aligned}
\langle c((g_j(\xi_j)),\tau\rangle
&=\frac{1}{2\pi\ii}\int_C\bar\partial h\wedge\tau\\
&=\frac{1}{2\pi\ii}\int_C\dd(h\tau)\\
&=\lim_{\eps\to0}\frac{1}{2\pi\ii}\int_{C_\eps}\dd(h\tau)\\
&=\lim_{\eps\to0}\frac{1}{2\pi\ii}\int_{\partial C_\eps}h\tau\\
&=\lim_{\eps\to0}\sum_{j=1}^N\frac{1}{2\pi\ii}\int_{\partial D_\eps^{(j)}}(h_j-s_j)\tau\\
&\qquad\text{（沿 }\alpha_i\text{、}\beta_i\text{、}l_j\text{ 的积分相消）}\\
&=\lim_{\eps\to0}\sum_{j=1}^N\frac{1}{2\pi\ii}\Bigl\{\int_{\partial D_\eps^{(j)}}h_j\tau-\int_{\partial D_\eps^{(j)}}s_j\tau\Bigr\}\\
&=\sum_{j=1}^N\operatorname*{Res}_{Q_j}(h_j\tau)-\lim_{\eps\to0}\frac{1}{2\pi\ii}\int_{D_\eps^{(j)}}\dd(s_j\tau) .
\end{aligned}
\]

% [Tutor Note: 笔误（两处）：(i) 上式第三行右端印作 $h_a-s_a=h_b-s_a$，但合拢上链需要
%  右边减 $s_b$；由 $h_a-h_b=s_b-s_a$（即 $h_{ab}=s_b-s_a$）得 $h_a-s_a=h_b-s_b$，
%  才与 $a,b$ 无关地定义 $h$。(ii) 最后一行的第二项印为 $\int_{D_\eps^{(j)}}$（无 $\sum_j$），
%  但它是从上一行 $\sum_{j=1}^N\frac{1}{2\pi\ii}\int_{\partial D_\eps^{(j)}}s_j\tau$ 逐项
%  用 Stokes 得到的，应带 $\sum_{j=1}^N$。两处均照原文排，不改。原文 p.24]
% [Tutor Note: 首行 $\int_C\bar\partial h\wedge\tau$ 与末行留数公式之间，原文省略了
%  $\tau$ 是 $(1,0)$-形式（否则 $\bar\partial h\wedge\tau$ 不是 $2$-形式、无法积分）；
%  这也说明 p.21 的配对里第二个因子应是 $\omega_C(m\sum Q_j)$ 而非 $\cO_C(m\sum Q_j)$。原文 p.24]

由于 $s_j\tau$ 在 $Q_j$ 的邻域内是 $C^\infty$ 的，我们有
\[
\lim_{\eps\to0}\int_{D_\eps^{(j)}}\dd(s_j\tau)=0 .
\]
此外我们有 $\operatorname{Res}_{Q_j}(h_j\tau)=\operatorname{Res}_{\xi_j=0}(g_j(\xi_j)\tau_j)$，其中 $\tau_j$ 是 $\tau$ 在 $Q_j$ 处关于 $\xi_j$ 的 Laurent 展开；于是我们断言
\[
\langle c((g_j(\xi_j)),\tau\rangle=\sum_{j=1}^N\operatorname*{Res}_{\xi_j=0}g_j(\xi_j)\tau_j .
\]
这就是所要的结果。

次设 $C$ 有结点。为简单起见，设 $C$ 只有一个结点 $P$。设 $\pi:\tilde C\to C$ 是正规化映射，$\pi^{-1}(P)=\{P_+,P_-\}$。则 $\pi^\ast\omega_C=\omega_{\tilde C}(P_++P_-)$。因此
\[
\pi^\ast H^0\Bigl(C,\omega_C\bigl(m\sum_{j=1}^N Q_j\bigr)\Bigr)\subset H^0\Bigl(\tilde C,\omega_{\tilde C}\bigl(m\sum_{j=1}^N Q_j\bigr)\Bigr).
\]

% [原文 p.25]
这说明在 $\tilde C$ 上 $(g_j(\xi_j))$ 被 $H^0(\tilde C,\omega_{\tilde C}(m\sum_{j=1}^N Q_j))$ 零化。于是由上面的论证，存在 $\tilde g\in H^0(\tilde C,\cO_{\tilde C}(m\sum_{j=1}^N Q_j))$ 使得 $p(\tilde g)=(g_j(\xi_j))$。另一方面 $(g_j(\xi_j))$ 被 $H^0(C,\omega_C(m\sum_{j=1}^N Q_j))$ 零化。因此对 $\tau\in\pi^\ast H^0(C,\omega_C(m\sum_{j=1}^N Q_j))\setminus H^0(\tilde C,\omega_{\tilde C}(m\sum_{j=1}^N Q_j))$ 我们有
\[
\sum_{j=1}^N\operatorname*{Res}_{Q_j}(\tilde g\tau)=0 .
\]
另一方面 $\tilde g\tau$ 在 $P_+$ 与 $P_-$ 处有一阶极点。因此我们断言
\begin{equation}
\operatorname*{Res}_{P_+}(\tilde g\tau)+\operatorname*{Res}_{P_-}(\tilde g\tau)=0 . \tag{1.52}
\end{equation}
由于我们有
\[
\operatorname*{Res}_{P_+}(\tilde g\tau)=\tilde g(P_+)\operatorname*{Res}_{P_+}\tau,\qquad
\operatorname*{Res}_{P_-}(\tilde g\tau)=\tilde g(P_-)\operatorname*{Res}_{P_-}\tau
\]
以及
\[
\operatorname*{Res}_{P_+}\tau+\operatorname*{Res}_{P_-}\tau=0 ,
\]
我们有
\[
\tilde g(P_+)=\tilde g(P_-) .
\]
这意味着 $\tilde g$ 是亚纯函数
\[
g\in H^0\Bigl(C,\cO_C\bigl(m\sum_{j=1}^N Q_j\bigr)\Bigr)
\]
的拉回。这说明 (1.51) 对结点曲线也成立。
关于 $H^0(C,\omega_C(\ast\sum_{j=1}^N Q_j))$ 的刻画，可以类似地证明。\hfill$\square$

% [Tutor Note: 关键跳步（★★）：由 $\tilde g(P_+)=\tilde g(P_-)$ 推出 $\tilde g$ 是 $C$ 上
%  亚纯函数的拉回。理由：$\pi$ 在 $\tilde C\setminus\{P_+,P_-\}$ 上双全纯到
%  $C\setminus\{P\}$，故 $\tilde g$ 在两侧的极限值相同即可沿 $\pi$ 下降（Riemann 可去
%  奇点定理的反向用法：$C$ 上定义 $g(x)=\tilde g(\pi^{-1}(x))$，在 $P$ 处由等值给出
%  单值极限）。条件 $\tilde g\in H^0(\cO_{\tilde C}(m\sum Q_j))$ 保证极点阶不增。原文 p.25]
% [Tutor Note: $\pi^\ast\omega_C=\omega_{\tilde C}(P_++P_-)$ 是结点曲线对偶化层的基本事实：
%  取 $\tilde C\setminus\{P_\pm\}$ 上的全纯 $1$-形式，在 $P_\pm$ 处允许一阶极点；
%  与 (1.46) 合起来看，$\omega_C$ 正是"两侧留数相等"的那一半，故拉回到 $\tilde C$ 后
%  两侧各有一阶极点。原文 p.25]
